% Generated draft: requires mathematical and source review.
\documentclass[11pt]{article}
\usepackage[margin=25mm]{geometry}
\usepackage{fontspec}
\setmainfont{Latin Modern Roman}
\usepackage{amsmath,amssymb,mathtools,unicode-math}
\setmathfont{Latin Modern Math}
\usepackage{enumitem,needspace,xurl,hyperref}
\setlength{\parindent}{0pt}
\setlength{\parskip}{5pt}
\setlength{\emergencystretch}{3em}
\newcommand{\N}{\mathbb{N}}
\newcommand{\Z}{\mathbb{Z}}
\newcommand{\Q}{\mathbb{Q}}
\newcommand{\R}{\mathbb{R}}
\newcommand{\C}{\mathbb{C}}
\newcommand{\F}{\mathbb{F}}
\newcommand{\A}{\mathbb{A}}
\newcommand{\PP}{\mathbb{P}}
\newcommand{\E}{\mathbb{E}}
\newcommand{\eps}{\varepsilon}
\newcommand{\dd}{\,\mathrm{d}}
\begin{document}
\section*{UnsolvedMath ID: 30002366}
\textbf{Source problem code:} catalogue entry
\subsection{Definitions and mathematical statement}
The supplied record asks for an unstable motivic homotopy exact sequence. Its base field and admissible field hypotheses are not specified. The following formulation uses the conventional Nisnevich setting to explain the notation; it does not assert that the conjecture is intended over every field.

Fix a base field \(k\), subject to hypotheses that remain to be recovered from the source, and let \(\mathrm{Sm}_k\) denote the category of smooth schemes of finite type over \(k\), equipped with the Nisnevich topology. An abelian sheaf \(F\) is strictly \(\mathbb A^1\)-invariant if, for every \(U\in\mathrm{Sm}_k\) and every integer \(i\geq 0\), projection induces an isomorphism
\[
H^i_{\mathrm{Nis}}(U,F)\xrightarrow{\sim}H^i_{\mathrm{Nis}}(U\times\mathbb A^1,F).
\]
Its contraction is the sheaf
\[
F_{-1}(U)=\ker\!\left(F(U\times\mathbb G_m)\xrightarrow{1^*}F(U)\right),
\]
where \(1^*\) is restriction along the unit section. Define \(F_{-r}=(F_{-(r-1)})_{-1}\) for \(r\geq 2\).

Write \(K^M_j\) for the unramified Milnor \(K\)-theory sheaf and \(K^M_j/24\) for the cokernel sheaf of multiplication by \(24\). Write \(\pi_n^{\mathbb A^1}(\mathbb A^n\setminus\{0\})\) for the \(n\)-th \(\mathbb A^1\)-homotopy sheaf of punctured affine space, conventionally pointed at \((1,0,\ldots,0)\). The symbol \(GW_{n+1}^{n}\) is retained from the clean statement and is interpreted as a higher Grothendieck--Witt sheaf, with lower index denoting degree and upper index denoting the duality shift. The intended indexing convention and its field hypotheses require source confirmation.

\textbf{Proposed statement, with unresolved base hypotheses.} For every integer \(n\geq 4\), do there exist morphisms of strictly \(\mathbb A^1\)-invariant sheaves
\[
K^M_{n+2}/24\xrightarrow{\alpha_n}
\pi_n^{\mathbb A^1}(\mathbb A^n\setminus\{0\})
\xrightarrow{\beta_n}GW_{n+1}^{n}
\]
that are exact at the middle term and whose \(n\)-fold contractions form a short exact sequence
\[
0\longrightarrow (K^M_{n+2}/24)_{-n}
\xrightarrow{(\alpha_n)_{-n}}
\bigl(\pi_n^{\mathbb A^1}(\mathbb A^n\setminus\{0\})\bigr)_{-n}
\xrightarrow{(\beta_n)_{-n}}
(GW_{n+1}^{n})_{-n}
\longrightarrow 0?
\]
Here exactness of the uncontracted displayed sequence is interpreted minimally as \(\operatorname{im}(\alpha_n)=\ker(\beta_n)\), since the supplied display has no endpoint zeros. Whether surjectivity of \(\beta_n\) before contraction is also intended remains unresolved. The record does not describe the intended maps.

The original statement instead uses \(K^{\mathbb A^1}_{n+2}/24\) and \(GW_{n+1}\). The display above follows the supplied clean statement, but the record alone does not establish that these original symbols are equivalent to its replacements. The missing base hypotheses, indexing conventions, and endpoint interpretation prevent an unconditional precise formulation.
\subsection{Short English statement}
For each integer \(n\geq 4\), the supplied conjecture asks for an exact sequence connecting Milnor \(K\)-theory modulo \(24\), an unstable motivic homotopy sheaf of punctured affine space, and a higher Grothendieck--Witt sheaf, becoming short exact after \(n\) contractions. The base field hypotheses, intended maps, and some notation remain unresolved. The supplied partially solved label is dated curation, not an independently verified mathematical conclusion.
\subsection{Sources}
\begin{itemize}
\item Local UnsolvedMath registry, supplied complete source record, permanent ID \(30002366\), title \emph{Unstable Motivic Homotopy Exact-Sequence Conjecture}, source problem code \texttt{catalogue entry}, in the Oberwolfach Reports Open Problems collection. No individual proposer is recorded; the supplied proposed year is \(2013\).
\item Supplied report citation: \emph{Algebraic K-theory and Motivic Cohomology} (2014), Oberwolfach Reports, pp.\ 1861--1913, DOI \texttt{10.4171/owr/2013/32}. \url{https://doi.org/10.4171/owr/2013/32}. This is the report provenance supplied by the record; the report itself has not been consulted for this draft.
\item Supplied literature source, ``Oberwolfach report on motivic splitting'': \url{https://ems.press/content/serial-article-files/46461?nt=1}.
\item Supplied literature source, ``Asok and Fasel, vector bundles'': \url{https://ems.press/content/book-chapter-files/33191?nt=1}.
\item Supplied literature source, ``Rondigs and Spitzweck, motivic homotopy'': \url{https://msp.org/obs/2025/6-1/obs-v6-n1-p06-s.pdf}.
\item Literature assessment dated exactly \texttt{2026-08-21}: the supplied curation labels the problem partially solved and reports that the stable \(2\)-primary statement over fields of \(2\)-cohomological dimension at most \(2\), together with several low-dimensional splitting cases, is proved, while the full unstable conjecture and its secondary obstruction remain open. This is supplied dated curation derived from the OWR collection, not independently verified here; its time-sensitive status and cited claims require independent checking.
\item The registry labels the statement \texttt{corrected\_verified} and reports that it was checked against the official Oberwolfach report, with hypotheses and notation restored. That verification claim is supplied metadata, not a verification performed for this draft. In particular, the supplied clean statement still does not state base field hypotheses.
\end{itemize}
\end{document}
